\documentclass[11pt]{article}
\PassOptionsToPackage{numbers,sort&compress}{natbib}
\usepackage[final]{acl}
\usepackage{times}
\usepackage{latexsym}
\usepackage[T1]{fontenc}
\usepackage[utf8]{inputenc}
\usepackage{microtype}

\usepackage{amsmath,amssymb}
\usepackage{booktabs,array}
\usepackage[table]{xcolor}
\usepackage{tikz}
\usetikzlibrary{arrows.meta,positioning}
\usepackage{float}
\usepackage{graphicx}

\definecolor{pipeblue}{RGB}{30,100,180}
\definecolor{pipeteal}{RGB}{0,140,130}
\definecolor{pipeorange}{RGB}{210,100,20}
\definecolor{pipegreen}{RGB}{30,140,60}
\definecolor{pipegray}{RGB}{80,80,90}
\definecolor{headerblue}{RGB}{220,235,255}
\definecolor{rowtint}{RGB}{245,250,255}

\setlength\titlebox{7\baselineskip}

\title{PrimeLine@DravidianLangTech 2026: Hope Speech Detection in Tulu\\
Using XLM-RoBERTa for Coarse and Fine-Grained Classification}

\author{Rithikaa~V \quad Sanjay~Krishnan~K \quad Nithya~Varshini~C~N~R \\
  \textit{Guided by: S.~Sumathi} \\
  Department of Information Technology, St.~Joseph's College of Engineering \\
  \url{[GitHub link to be added]}
}

\begin{document}
\maketitle

\begin{abstract}
Hope speech detection in low-resource, code-mixed languages presents a genuine
challenge for natural language processing. Tulu, a Dravidian language spoken along
the coastal regions of Karnataka and Kerala, is one such language where social media
content is deeply code-mixed, blending Tulu, Kannada script, and English within a
single comment. Two classification tasks are addressed: a four-class coarse-grained
setting (Track~1) and a five-class fine-grained setting (Track~2). XLM-RoBERTa,
a cross-lingual transformer pre-trained on more than 100 languages, is fine-tuned
on the task-provided datasets using Google Colab with an NVIDIA T4 GPU. The system
achieves a Macro F1-score of 0.34 on Track~1 and 0.19 on Track~2 on the official
Codabench evaluation, establishing the first transformer-based baseline for hope
speech classification in Tulu.
\end{abstract}

\section{Introduction}

Hope speech refers to positive, forward-looking online content that encourages,
supports, or inspires others. As NLP research has expanded its scope, hope speech
detection has grown into a meaningful area of study, offering a constructive
counterpart to the widely explored tasks of hate speech and toxicity detection.
While considerable effort has been directed at identifying harmful content, the
computational identification of positive and hopeful discourse has received
comparatively less attention.

Identifying hopeful language in regional and minority languages is particularly
valuable, as it gives visibility to communities rarely represented in large-scale
language studies. For speakers of languages such as Tulu, computational tools can
open up possibilities for social listening and content moderation tailored to their
specific needs.

Tulu is spoken by roughly four to five million people in the Tulu Nadu region of
coastal Karnataka and Kerala. Despite its significant speaker population, Tulu remains
almost entirely absent from NLP research. There is no dedicated pre-trained language
model for Tulu and no established NLP pipeline for the language. Its presence on
social media is marked by heavy code-mixing, where native Tulu, Kannada script,
English, and emoji are blended freely within a single comment, making automated
classification considerably harder than equivalent tasks in monolingual settings.

The DravidianLangTech shared task on Hope Speech Detection in Tulu,
organized as part of ACL 2026, provides the first structured benchmark for this
problem. Track~1 covers four coarse-grained categories: Encouraging Hope,
Discouraging Hope, Blended Hope, and Uninvolved. Track~2 covers five fine-grained
categories: Inspiring Hope, Hopelessness, Realistic Hope, Optimistic Hope, and
Fading Hope. The transition to fine-grained classification makes the task
progressively more demanding, as Track~2 categories differ in degree rather than
in kind, with genuinely ambiguous boundaries.

The approach taken in this work centers on fine-tuning XLM-RoBERTa \citep{conneau2020},
a multilingual transformer pre-trained on more than 100 languages. Its cross-lingual
representations allow it to handle Tulu\textquotesingle{}s mixed scripts without any
language identification or transliteration steps, making it a practical and principled
choice for this low-resource, code-mixed setting. The system is trained end-to-end
with a linear classification head on top of the encoder, and predictions are generated
by taking the argmax of the softmax output.

\section{Related Work}

Hope speech detection as a formal NLP task was first introduced by
\newcite{chakravarthi2020} through the HopeEDI dataset, covering English, Tamil, and
Malayalam. The dataset and baselines demonstrated that transformer-based models
clearly outperform classical approaches even on small multilingual corpora, setting
the stage for the shared tasks that followed.

\newcite{chakravarthi2021findings} reported findings from a shared task on hope speech
detection across English, Tamil, and Malayalam, confirming that multilingual
transformers consistently outperform classical machine learning approaches on
code-mixed text.

Several systems from those shared tasks demonstrated the strength of XLM-RoBERTa.
\newcite{ziehe2021} showed that fine-tuning XLM-RoBERTa achieved strong results
across English, Malayalam, and Tamil without any task-specific feature engineering.
\newcite{mahajan2021} ranked first in both English and Tamil by fine-tuning RoBERTa
and XLM-RoBERTa respectively. \newcite{balouchzahi2021} built an XLM-RoBERTa
ensemble specifically for code-mixed Dravidian text. \newcite{huang2021} combined
XLM-RoBERTa with TF-IDF features and ranked first in both English and Malayalam.
\newcite{singh2021} explored multilingual transformer models pre-trained on Indian
languages and achieved second place in both English and Malayalam. These results
collectively established XLM-RoBERTa as the dominant baseline for Dravidian hope
speech classification.

On the model side, \newcite{vaswani2017} introduced the Transformer architecture,
which set the foundation for the modern pre-trained language model paradigm.
\newcite{devlin2019} built on this with BERT, demonstrating the power of masked
language modelling for downstream NLP tasks. \newcite{liu2019} followed with RoBERTa,
showing that more robust pre-training strategies yield better performance.
\newcite{conneau2020} extended this to the multilingual setting with XLM-RoBERTa,
pre-trained on over 100 languages including several Indic scripts.

Dravidian language NLP has received growing attention in recent years.
\newcite{chakravarthi2022overview} reported a comprehensive shared task overview on
hope speech across five languages. \newcite{chakravarthi2022multilingual} further
presented a multilingual hope speech dataset covering English and Dravidian languages.
\newcite{sundar2022} proposed a stacked encoder architecture using cross-lingual
embeddings and demonstrated improvements over standard fine-tuning baselines.

Research on Tulu specifically is extremely sparse. \newcite{shetty2023} examined
sentiment analysis in code-mixed Tulu and Tamil text, noting the severe scarcity of
annotated data. \newcite{ehsan2023} tackled Tulu sentiment classification using BiLSTM
models with ELMo embeddings, identifying the lack of Tulu-specific language resources
as the primary bottleneck. Neither work addresses hope speech detection in Tulu;
the current work is the first to do so.

\section{Dataset and Preprocessing}

\subsection{Dataset}
The datasets for both tracks were prepared by the shared task organizers and drawn
from social media platforms used by Tulu speakers. The text reflects real-world Tulu
online writing, mixing Tulu, Kannada script, English, and emoji freely within
individual comments. This code-mixed, multi-script nature is a defining feature of
the data and makes classification non-trivial even for human annotators.

Track~1 contains 5,991 training samples across four coarse categories. The
distribution is imbalanced, with Uninvolved accounting for 2,490 samples (41.6\%)
and Discouraging Hope having the fewest at 711 (11.9\%). Track~2 has 3,185 training
samples across five fine-grained labels, with Inspiring Hope as the majority class
at 35.4\% and Fading Hope as the minority at 7.4\%.
Tables~\ref{tab:track1} and~\ref{tab:track2} show the full distributions.

\begin{table}[H]
\centering\small
\setlength{\tabcolsep}{5pt}\renewcommand{\arraystretch}{1.1}
\rowcolors{2}{rowtint}{white}
\begin{tabular}{lcc}
\toprule
\rowcolor{headerblue}\textbf{Label} & \textbf{Count} & \textbf{\%} \\
\midrule
Encouraging Hope  & 1895 & 31.6\% \\
Discouraging Hope & 711  & 11.9\% \\
Blended Hope      & 895  & 14.9\% \\
Uninvolved        & 2490 & 41.6\% \\
\midrule
\textbf{Total} & \textbf{5991} & \textbf{100\%} \\
\bottomrule
\end{tabular}
\caption{Class distribution -- Track~1 (Coarse-Grained).}
\label{tab:track1}
\end{table}

\begin{table}[H]
\centering\small
\setlength{\tabcolsep}{4pt}\renewcommand{\arraystretch}{1.1}
\rowcolors{2}{rowtint}{white}
\begin{tabular}{lcc}
\toprule
\rowcolor{headerblue}\textbf{Label} & \textbf{Count} & \textbf{\%} \\
\midrule
Inspiring Hope  & 1129 & 35.4\% \\
Hopelessness    & 937  & 29.4\% \\
Realistic Hope  & 503  & 15.8\% \\
Optimistic Hope & 380  & 11.9\% \\
Fading Hope     & 236  &  7.4\% \\
\midrule
\textbf{Total} & \textbf{3185} & \textbf{100\%} \\
\bottomrule
\end{tabular}
\caption{Class distribution -- Track~2 (Fine-Grained).}
\label{tab:track2}
\end{table}

\subsection{Preprocessing}
Preprocessing was kept minimal by design. Text samples were passed directly to the
XLM-RoBERTa tokenizer, which handles diverse scripts through a SentencePiece-based
subword vocabulary. Padding and truncation were applied at a maximum length of 128
tokens. Transliteration, language identification, and stopword removal were
deliberately skipped, as these steps risk erasing code-switching cues that carry
genuine meaning. No normalization of Kannada or Roman script was applied, preserving
the natural variation present in user-generated content. This minimal preprocessing
strategy is consistent with prior work on code-mixed Dravidian NLP, where aggressive
preprocessing has been shown to hurt rather than help performance.

\section{Methodology}

\subsection{System Pipeline}
The system follows a standard transformer fine-tuning pipeline. Each raw text sample
is first passed to the XLM-RoBERTa tokenizer, which applies SentencePiece subword
segmentation and encodes the input as a sequence of token IDs with a special
\texttt{[CLS]} token prepended. The encoded sequence is fed into the XLM-RoBERTa
encoder, which produces contextualized representations through twelve layers of
multi-headed self-attention. The hidden state of the \texttt{[CLS]} token at the
final encoder layer is extracted and passed through a linear classification head.
Output logits are converted to class probabilities via softmax, and the predicted
label is the class with the highest probability. For Track~1, the output layer has
four neurons; for Track~2, five. Figure~\ref{fig:pipeline} illustrates the pipeline.

\begin{figure}[H]
\centering
\begin{tikzpicture}[node distance=3mm,
  box/.style={rectangle,rounded corners=2pt,minimum width=\columnwidth-2mm,
    minimum height=7mm, align=center,font=\small\bfseries,draw,thick},
  arr/.style={-Stealth,thick,color=pipegray}]
\node[box,fill=pipeorange!25,draw=pipeorange!80!black]
  (inp) {Raw Text Input (Code-Mixed Tulu)};
\node[box,fill=pipeteal!20,draw=pipeteal!80!black,below=of inp]
  (tok) {XLM-RoBERTa Tokenizer (SentencePiece, max 128)};
\node[box,fill=pipeblue!20,draw=pipeblue!80!black,below=of tok]
  (enc) {XLM-RoBERTa Encoder (12 layers, 768 dims)};
\node[box,fill=pipeblue!10,draw=pipeblue!60!black,below=of enc]
  (cls) {\texttt{[CLS]} Token Representation};
\node[box,fill=pipegreen!20,draw=pipegreen!80!black,below=of cls]
  (hd) {Linear Head: 4 classes (T1) / 5 classes (T2)};
\node[box,fill=pipeorange!30,draw=pipeorange!80!black,below=of hd]
  (out) {Predicted Hope Speech Label};
\draw[arr](inp)--(tok);\draw[arr](tok)--(enc);\draw[arr](enc)--(cls);
\draw[arr](cls)--(hd);\draw[arr](hd)--(out);
\end{tikzpicture}
\caption{End-to-end system pipeline for hope speech classification in Tulu.}
\label{fig:pipeline}
\end{figure}

\subsection{Model Formulation}
The core of the system is \texttt{xlm-roberta-base}, a transformer encoder built on
the architecture of \newcite{vaswani2017} and trained using the robust pre-training
strategy of \newcite{liu2019} extended to 100+ languages by \newcite{conneau2020}.
The model contains approximately 270 million parameters and was pre-trained on
filtered Common Crawl data spanning 100 languages, including Indic scripts.

Within each transformer layer, the self-attention mechanism computes attention
scores over all token pairs. Given query matrix $Q$, key matrix $K$, and value
matrix $V$, scaled dot-product attention is:
\begin{equation}
\text{Attention}(Q,K,V) = \text{softmax}\!\left(\frac{QK^{\top}}{\sqrt{d_k}}\right)V
\end{equation}
where $d_k$ is the key vector dimensionality. This allows the model to capture
long-range dependencies across mixed scripts and languages simultaneously.

The \texttt{[CLS]} representation is fed into a linear head, and logits
$z_1,\ldots,z_C$ are normalized via softmax to produce class probabilities:
\begin{equation}
P(y = c \mid \mathbf{x}) = \frac{\exp(z_c)}{\sum_{j=1}^{C}\exp(z_j)}
\end{equation}

The model is trained end-to-end by minimizing categorical cross-entropy loss:
\begin{equation}
\mathcal{L} = -\frac{1}{N}\sum_{i=1}^{N}\sum_{c=1}^{C}
  y_{ic}\log P(y = c \mid \mathbf{x}_i)
\end{equation}
where $N$ is the number of training samples, $C$ the number of classes, and
$y_{ic}\in\{0,1\}$ the one-hot ground truth indicator. Gradients are
back-propagated through the full encoder, adapting pre-trained representations
to the specific vocabulary and style of Tulu social media text.

\subsection{Training Setup}
Both models were fine-tuned on Google Colab with an NVIDIA T4 GPU for 2~epochs
with batch size~8 and AdamW optimizer (learning rate $2{\times}10^{-5}$, weight
decay $0.01$). A linear warmup was applied over the first 6\% of training steps,
followed by linear decay to zero. The development set served as the validation split
throughout, and the best checkpoint by validation loss was selected for inference.
Predictions on test sets were obtained by argmax over output logits, converted to
label strings, and submitted to Codabench. No ensembling was applied.

\section{Results and Discussion}

\subsection{Official Results}
Table~\ref{tab:results} reports official Codabench evaluation scores. Accuracy and
macro-averaged precision, recall, and F1 are reported. Macro-averaging is the primary
metric as it weights all classes equally, making it sensitive to performance on
minority categories regardless of their frequency.

\begin{table}[H]
\centering\small
\setlength{\tabcolsep}{4pt}\renewcommand{\arraystretch}{1.1}
\rowcolors{2}{rowtint}{white}
\begin{tabular}{lcccc}
\toprule
\rowcolor{headerblue}\textbf{Track} & \textbf{Acc} & \textbf{P} & \textbf{R} & \textbf{F1} \\
\midrule
Track 1 -- Coarse & 0.60 & 0.30 & 0.41 & 0.34 \\
Track 2 -- Fine   & 0.24 & 0.19 & 0.19 & 0.19 \\
\bottomrule
\end{tabular}
\caption{Macro-averaged results on Codabench (P\,=\,Precision, R\,=\,Recall).}
\label{tab:results}
\end{table}

On Track~1, an accuracy of 0.60 is achieved with a Macro F1 of 0.34. The gap
between these two figures reflects a model that performs well on the majority
Uninvolved class but struggles with less frequent categories. Track~2 is harder,
with accuracy at 0.24 and Macro F1 at 0.19, reflecting the semantic closeness of
its five fine-grained categories and severe class imbalance.

\subsection{Error Analysis}
A detailed look at predictions reveals that the accuracy-F1 gap on Track~1 is driven
primarily by over-prediction of the Uninvolved class (41.6\% of training data).
Samples from Blended Hope and Discouraging Hope are prone to being absorbed into
this dominant category, as they share surface-level vocabulary with neutral comments.
The decision boundary between Encouraging Hope and Blended Hope is especially
difficult, since a comment can simultaneously express genuine positivity and realistic
acknowledgment of hardship, making clean categorical assignment ambiguous.

Track~2 exhibits a more extreme version of the same problem. The five fine-grained
categories are semantically adjacent and differ in degree rather than kind. A comment
expressing uncertainty about the future could plausibly belong to Fading Hope,
Hopelessness, or Realistic Hope depending on subtle contextual cues difficult to
capture with only two epochs of fine-tuning. The scarcity of Fading Hope samples
(236 of 3,185) further limits learning for this class.

\subsection{Discussion and Future Directions}
The results demonstrate that XLM-RoBERTa provides a viable starting point for
hope speech detection in Tulu even without Tulu-specific pre-training. Future work
should address class imbalance via weighted loss functions or back-translation
augmentation from Kannada. Training for more epochs with early stopping based on
macro F1 could help balance performance across classes. Larger multilingual models
or Dravidian-family-specific pre-trained models could yield better representations
for Tulu\textquotesingle{}s lexical patterns. Contrastive learning objectives that
explicitly separate semantically similar hope categories could also resolve the
fine-grained boundary confusion observed in Track~2.

\section{Conclusion}
This work presents the first transformer-based system for hope speech detection in
code-mixed Tulu, fine-tuned on datasets from the DravidianLangTech@ACL~2026 shared
task. XLM-RoBERTa is applied to both tracks, achieving Macro F1 scores of 0.34 and
0.19 on Track~1 and Track~2 respectively. The results establish a meaningful baseline
for future work and highlight the dual challenge of class imbalance and fine-grained
semantic similarity. Despite modest absolute scores, the cross-lingual transfer
capabilities of XLM-RoBERTa prove sufficient to produce non-trivial predictions
even for a language with no dedicated NLP resources. It is hoped that this
contribution encourages further computational research on Tulu and other low-resource
Dravidian languages that remain largely absent from the NLP literature.

\section*{Limitations}
The primary limitation is the use of only two training epochs, restricting the
model's ability to learn reliable decision boundaries for minority classes. The
absence of Tulu-specific pre-trained representations means cross-lingual transfer
may not fully capture the language's phonological and morphological patterns. Class
imbalance was not addressed through resampling or loss-weighting, suppressing
performance on rare categories. The system does not perform language identification
or script normalization, both avenues for future improvement.

\section*{Acknowledgments}
The organizers of the DravidianLangTech@ACL~2026 shared task are thanked for
providing the datasets and the Codabench evaluation platform. Sincere gratitude is
expressed to the project guide, S.~Sumathi, Department of Information Technology,
St.~Joseph's College of Engineering, for constant guidance throughout this work.

\bibliography{custom}

\end{document}